\chapter{L'état de l'art}
\label{chap:etat}

\begin{chapeau}
Ce chapitre résume ce qui est publié et utile au projet : la méthode CORA, les valeurs normales du chien, la mesure de
la torsion, les guides imprimés, la segmentation des os et l'ostéotomie en dôme. Il se termine par la liste des
décisions que nous en avons tirées.
\end{chapeau}

\section{La méthode CORA de Paley}
\label{sec:cora-paley}

\subsection{Le principe}

La méthode de Paley \cite{paley2002} part d'une idée simple : un os déformé à un seul apex est fait de deux segments
droits. Chaque segment a son axe. On construit chaque axe à partir de la ligne articulaire de son extrémité et de
l'angle articulaire \emph{normal} : par exemple, l'axe du segment proximal du tibia est la droite qui fait l'angle
mMPTA normal (environ 93\degre{}) avec la ligne du plateau tibial. Les deux axes se croisent en un point : le
\textbf{CORA}, pour \emph{center of rotation of angulation}. L'angle entre les axes est l'angle de la déformité
$\theta$ (\cref{fig:methode-cora}).

\fig[0.8]{methode_cora.png}{Principe de la méthode CORA : (a) les axes des deux segments se croisent au CORA et font
l'angle $\theta$ ; (b) en coupant au CORA et en redressant autour de ce point, les deux axes sont de nouveau
alignés.}{fig:methode-cora}

\subsection{L'axe de correction et la bissectrice}

Redresser l'os, c'est faire tourner le segment distal de l'angle $\theta$ autour d'un axe perpendiculaire au plan de
la déformité : l'\textbf{axe de correction de l'angulation}, ou ACA. Paley montre que l'ACA doit se trouver sur la
\textbf{bissectrice transverse} du CORA, la droite qui partage en deux l'angle entre les axes, dans le plan de la
déformité. C'est la seule façon de réaligner les axes sans translation. La position de l'ACA le long de la bissectrice
choisit le type de coin :
\begin{itemize}
  \item ACA sur la corticale concave : \textbf{coin de fermeture} (l'os raccourcit) ;
  \item ACA sur la corticale convexe : \textbf{coin d'ouverture} (l'os s'allonge) ;
  \item ACA au milieu de l'os : \textbf{coin neutre} (une moitié se ferme, l'autre s'ouvre ; la longueur change peu).
\end{itemize}

\subsection{Les règles d'ostéotomie}

Paley énonce trois règles, selon l'endroit de la coupe par rapport au CORA :
\begin{itemize}
  \item \textbf{règle 1} : la coupe et l'ACA passent par le CORA. Les axes sont réalignés, sans translation ;
  \item \textbf{règle 2} : la coupe est faite à un autre niveau, mais l'ACA reste au CORA. Les axes sont réalignés,
  et le fragment distal glisse le long de la coupe : c'est une \emph{translation à l'ostéotomie} ;
  \item \textbf{règle 3} : la coupe et l'ACA sont loin du CORA. L'angle est corrigé, mais les axes restent décalés :
  il reste une translation de l'axe, à éviter.
\end{itemize}
On utilise la règle 2 quand on ne peut pas couper au CORA, par exemple parce qu'il est trop près de l'articulation.
Notre logiciel applique la règle 1 par défaut et passe à la règle 2 si le trait de coupe serait trop près du
cartilage (\cref{sec:regle2}).

\subsection{Le plan oblique et l'angle vrai}
\label{sec:oblique}

Une déformité réelle est rarement pure. Le plus souvent, elle a une composante de face $F$ et une composante de profil
$S$ : elle se trouve dans un plan \emph{oblique}. L'angle vrai n'est pas la somme des deux, et corriger les deux vues
l'une après l'autre reviendrait à mal corriger. Quand l'axe proximal est vertical, Paley donne
\begin{equation}
  \tan^2\theta = \tan^2 F + \tan^2 S,
  \qquad
  \tan\varphi = \frac{\abs{\tan S}}{\abs{\tan F}},
  \label{eq:oblique}
\end{equation}
où $\varphi$ est l'orientation du plan de la déformité par rapport au plan frontal \cite{paley2002,qatu2024}. Pour
$F = 15\degre$ et $S = 10\degre$, on trouve $\theta = 17{,}8\degre$, et non 25\degre. Une seule ostéotomie oblique
corrige alors les deux composantes à la fois \cite{kim2017}. La démonstration est donnée à la \cref{sec:angle-vrai}.

\section{Les valeurs normales chez le chien}
\label{sec:normales}

Les valeurs normales des angles articulaires viennent de radiographies de chiens sains, surtout de grandes races
(\cref{tab:normales}). Deux points comptent :
\begin{itemize}
  \item \textbf{la race compte}. L'aLDFA normal va de 94 à 98\degre{} selon la race \cite{tomlinson2007} ; l'équipe
  Vezzoni retient 94 ± 3,3\degre{} (berger allemand) comme borne de normalité \cite{vezzoni2022}. Une étude au
  scanner compare six races \cite{bruhschwein2023r} ;
  \item \textbf{la définition compte}. Ces valeurs sont mesurées sur des radiographies, avec des repères placés par un
  observateur. Un logiciel qui mesure en 3D avec ses propres repères ne trouve pas exactement le même angle sur le même
  os. Nous corrigeons cet écart de définition (\cref{sec:calibration}).
\end{itemize}

\begin{table}[htbp]
\centering\small
\caption{Valeurs normales publiées (moyenne ± écart-type) utilisées par le logiciel. En gras : la référence par
défaut (travaux de l'équipe Vezzoni, sinon la littérature).}
\label{tab:normales}
\begin{tabularx}{\linewidth}{lL{2.6cm}Y}
\toprule
\titrecol{Angle} & \titrecol{Valeur} & \titrecol{Source} \\
\midrule
aLPFA & 103 ± 6,4\degre & Tomlinson 2007, Labrador \cite{tomlinson2007} \\
aLDFA & \textbf{94 ± 3,3\degre} & Vezzoni 2022 (valeur du berger allemand de Tomlinson 2007) \cite{vezzoni2022} \\
 & 97 ± 3,2\degre & Tomlinson 2007, Labrador \cite{tomlinson2007} \\
aCdDFA & 90,5 ± 6,3\degre & Kara 2018, scanner 3D, définition voisine \cite{kara2018} \\
mMPTA & \textbf{93,3 ± 1,78\degre} & Olsen, Vezzoni 2016 (= Dismukes 2007) \cite{olsen2016,dismukes2007} \\
mMDTA & 96,0 ± 2,7\degre & Dismukes 2007 \cite{dismukes2007} \\
mCaPTA & \textbf{63,7 ± 2,6\degre} & Guénégo, Vezzoni 2021 : 90\degre{} − pente du plateau (26,3\degre) \cite{guenego2021} \\
mCrDTA & 81,6 ± 3,2\degre & Dismukes 2008, radiographie de profil \cite{dismukes2008s} \\
valgus tibial & 11,8 ± 3,1\degre & cité par IMAIOS \cite{imaios} \\
pente du plateau (TPA) & 22,5\degre & lévrier sain, « à considérer comme normal » \cite{wilke2002} \\
\bottomrule
\end{tabularx}
\end{table}

\section{Mesurer la torsion}
\label{sec:etat-torsion}

La torsion est l'angle, autour de l'axe long de l'os, entre un axe transverse proximal et un axe transverse distal.
Au fémur, on parle d'\emph{antéversion} : l'angle entre l'axe du col et la ligne qui passe derrière les condyles.

\begin{itemize}
  \item \textbf{La radiographie ne suffit pas}. Elle ne distingue pas une torsion de l'os d'une simple rotation du
  membre sur la table \cite{apelt2005,aper2005}, et une rotation du fémur change à elle seule l'aLDFA mesuré
  \cite{miles2016}.
  \item \textbf{Le scanner mesure la torsion} \cite{dudley2006,mostafa2014}. Al Aiyan et al. trouvent une
  antéversion de 31 ± 5\degre{} chez 58 chiens de races moyennes et grandes \cite{alaiyan2021}. Brühschwein et al.
  proposent un repère 3D centré sur l'os, exact sur un modèle de torsion \cite{bruhschwein2023f}, et l'appliquent au
  tibia \cite{bruhschwein2023t}.
  \item \textbf{Seuils} : Longo et al. planifient des ostéotomies de dérotation en 3D et donnent des seuils de torsion
  anormale \cite{longo2022}.
  \item \textbf{Au bloc}, contrôler la rotation à l'œil ou à la radioscopie est imprécis : 4 à 9\degre{} d'erreur
  moyenne dans une étude sur cadavre \cite{ivanov2024} ; les signes radiographiques classiques ont leurs limites
  \cite{krettek1998,fang2015}.
\end{itemize}

\begin{attention}[title={Les conventions de signe varient}]
Selon les articles, une même torsion est dite « interne » ou « externe », et son signe change. Le logiciel suit les
Méthodes de Longo 2022 pour le tibia ($-$ = interne, $+$ = externe ; normale vers $-4\degre$, anormale en dessous de
$-10\degre$ ou au-dessus de $+2\degre$) et l'antéversion d'Al Aiyan 2021 pour le fémur ($+$ = col dirigé vers
crânial ; normale 31 ± 5\degre). Le résumé de l'article de Longo inverse les deux mots : il faut lire les Méthodes.
\end{attention}

\section{Planification virtuelle et guides imprimés}
\label{sec:etat-guides}

Plusieurs équipes ont publié des corrections planifiées sur ordinateur et réalisées avec des guides imprimés en 3D
(\cref{tab:guides}). Ces résultats portent sur toute la chaîne : préparation, guide et geste chirurgical.

\begin{table}[htbp]
\centering\small
\caption{Guides imprimés chez le chien : ce qui est publié.}
\label{tab:guides}
\begin{tabularx}{\linewidth}{L{2.6cm}L{3.2cm}Y}
\toprule
\titrecol{Étude} & \titrecol{Os, système} & \titrecol{Résultat} \\
\midrule
Worth 2019 \cite{worth2019} & avant-bras ; guides de coupe seuls & 1 membre sur 7 réopéré pour une rotation mal corrigée \\
Hall 2019 \cite{hall2019} & fémur distal (luxation de rotule) ; coupe + réduction & torsion à 1,7 ± 2,1\degre{} de la cible \\
Carwardine 2021 \cite{carwardine2021} & avant-bras ; coupe, repositionnement, plaques en titane & plaque sur mesure, trous percés à travers le guide \\
De Armond 2022 \cite{dearmond2022} & avant-bras, 1 ou 2 apex ; coupe + réduction & plus de 90~\% des paramètres à moins de 5\degre{} et 5~mm \\
Jeon 2023 \cite{jeon2023} & tibia de petites races & erreurs de 1,8 (profil), 2,3 (face) et 2,6\degre{} (torsion) \\
Bright 2024 \cite{bright2024} & tibia du teckel (pes varus) & moins de 2\degre{} et 1~mm dans tous les plans \\
\bottomrule
\end{tabularx}
\end{table}

Deux leçons en ressortent. D'abord, un guide de coupe seul contrôle mal la rotation ; il faut un moyen de
\textbf{remettre les fragments à la bonne place} après la coupe : un guide de réduction qui s'enfile sur des broches,
ou une plaque sur mesure. Ensuite, l'erreur finale est de l'ordre de \textbf{2\degre} : c'est l'objectif à viser, et
l'erreur du calcul doit être bien plus petite.

\section{Le choix du coin dans la pratique}
\label{sec:etat-coin}

Les publications donnent des stratégies types :
\begin{itemize}
  \item \textbf{pes varus du teckel} : coin d'ouverture médial au CORA, en gardant la corticale latérale comme
  charnière \cite{petazzoni2012} ;
  \item \textbf{varus du fémur distal} (luxation médiale de la rotule) : coin de fermeture latéral ; l'aLDFA passe de
  107,6 à 94,1\degre{} \cite{brower2017} ;
  \item \textbf{valgus du tibia proximal} : coin de fermeture médial associé à une TPLO ; le mMPTA passe de 101 à
  92,8\degre{} \cite{vezzoni2020} ;
  \item \textbf{pente du plateau tibial} : coin de fermeture crânial (CCWO), dont l'angle est relié à la pente
  \cite{guenego2021,guenego2025}. Un coin d'ouverture modifie facilement la pente : 1~mm d'erreur d'ouverture en avant
  change la pente d'environ 2\degre{} \cite{noyes2005,ducat2012}.
\end{itemize}

\section{La segmentation des os au scanner}
\label{sec:etat-segmentation}

Segmenter, c'est décider pour chaque voxel à quel os il appartient. Trois familles de méthodes sont utiles ici :
\begin{itemize}
  \item le \textbf{seuillage} : on garde les voxels au-dessus d'une valeur. Le seuil « Bone (CT) » du logiciel
  Mimics, 226~HU, est très utilisé en planification 3D. Mais à ce seuil, les os d'une articulation se touchent (volume
  partiel) et forment un seul bloc ;
  \item la \textbf{croissance de région} : on part d'un point (une graine) et on ajoute les voxels voisins qui
  ressemblent ;
  \item la \textbf{ligne de partage des eaux} (\emph{watershed}) \cite{meyer1994} : on voit l'image comme un relief
  et on le remplit d'eau depuis des sources ; les frontières sont là où deux eaux se rencontrent.
\end{itemize}
Le logiciel libre 3D Slicer \cite{fedorov2012} offre ces outils à la main. La surface 3D est ensuite extraite par
l'algorithme des \emph{marching cubes} \cite{lorensen1987}. Notre méthode (\cref{chap:segmentation}) combine un seuil
qui monte pour trouver une graine par os, puis une croissance de région compétitive.

\section{L'ostéotomie en dôme}
\label{sec:dome}

Quand on ne peut pas couper au CORA et qu'on veut éviter la translation de la règle 2, on peut faire une coupe
courbe, un cylindre centré sur le CORA : les fragments tournent l'un sur l'autre le long de cette coupe, sans
translation, en restant en contact. Elle est décrite au radius du chien, pour une ou deux déformités
\cite{fox2006}.

\begin{pourquoi}
Nous ne proposons pas le dôme : (1) il faut une lame courbe de rayon fixe et un guide qui fasse tourner la lame autour
d'un pivot, plus difficile à imprimer et à contrôler qu'une fente plane ; (2) un dôme cylindrique ne tourne qu'autour
de son axe : il corrige l'angle mais pas la torsion ; (3) au fémur distal et au tibia, le CORA est le plus souvent dans
une zone où la règle 1 est possible.
\end{pourquoi}

\section{Ce que nous en retenons}
\label{sec:retenu}

\begin{aretenir}[title={Les décisions du projet, tirées de l'état de l'art}]
\begin{enumerate}
  \item Partir du \textbf{scanner}, pour avoir la forme 3D et la torsion (\cref{chap:lecture}).
  \item Trouver les \textbf{repères automatiquement}, avec les définitions des articles, et corriger l'écart entre
  la mesure 2D publiée et notre mesure 3D (\cref{chap:reperes}).
  \item Faire la méthode CORA \textbf{en 3D} : un seul coin oblique, l'ACA sur la bissectrice, la règle 1 par défaut
  (\cref{chap:planification}).
  \item Corriger la \textbf{torsion} par une dérotation, avec les définitions publiées au scanner.
  \item Choisir le coin et la voie d'abord par des \textbf{règles explicites}, conformes aux pratiques publiées
  (\cref{sec:strategie}).
  \item Ne pas se contenter d'un guide de coupe : prévoir un \textbf{moyen de réduction} qui contrôle la rotation
  (guide de réduction, \cref{chap:guides}).
  \item Viser une erreur de calcul très inférieure aux \textbf{2\degre} publiés, et le vérifier sur des cas simulés
  (\cref{chap:validation}).
\end{enumerate}
\end{aretenir}
